\begin{tresc}
W turnieju szachowym brało udział 5 zawodników. Każdy z szachistów rozegrał jedną partię z każdym z pozostałych. Jeśli najlepszy zawodnik wygrał wszystkie partie, najsłabszy – przegrał wszystkie partie i odnotowano dwa remisy, to ile punktów uzyskał szachista, który zajął trzecie miejsce? Szachista otrzymywał jeden punkt za każde zwycięstwo, pół punktu za każdy remis i zero punktów za przegraną.
\end{tresc}
\begin{wariant}{A} 1 \end{wariant}
\begin{wariant}{B} 2 \end{wariant}
\begin{wariant}{C} 3 \end{wariant}
\begin{wariant}{D} 4 \end{wariant}
\begin{wariant}{E} 5 \end{wariant}
